\section{Machines: the operational carrier}
\label{sec:machines}

We now formalize the operational interface that the event-graph machine
instantiates.

\subsection{The generic machine}

The generic \(\mathsf{Machine}\) type is intentionally not specific to
VegasCore.  It contains a state type, per-player action types, internal
event types, public and private observation types, an outcome type,
initial state, availability predicates, total transition kernels,
terminality, observation functions, a partial outcome projection, and
utility.  Total transitions make trace composition uniform; legality is
kept as a separate predicate.

Here a machine is a precise operational model: what states exist,
which events are allowed, how events change state, what each player can
observe, and what payoff is read at the end.  The same interface can
describe the primitive event graph, an audited toy runtime, or a future
contract-level backend.

The outcome projection is partial.  This is important: nonterminal or
malformed states are not silently assigned zero utility.  For
compiler-generated reachable terminal states, payoff totality is proved
once at the graph-machine boundary.  Strategic kernel games are then
built from bounded runs whose terminal payoff projection is the checked
source payoff.

\subsection{Primitive event-graph machine}
\label{sec:mach:primitive}

The primitive graph machine interprets a compiled event graph one
available primitive event at a time.  Its state is a reachable graph
configuration, not an arbitrary raw configuration.  Player actions are
typed commit actions naming a graph node and value.  Internal events are
sample or reveal nodes.

Unavailable primitive events stutter in the total transition function,
but legal event laws and strategic presentations quantify over
availability evidence.  The stutter branch is only a totalization device:
the primitive machine does not choose which available event runs next.  A
scheduler, trace law, checkpoint policy, or runtime refinement supplies
that choice externally.  For example, a contract-execution machine might
execute many low-level steps before one protocol event is visible at the
graph boundary; a refinement proof later explains how that low-level
trace projects back to the primitive graph machine.

Public observations contain the completed-node downset and public
available field values.  Player observations contain the same
completed-downset graph information plus fields visible to that player
and ready commit information local to the player.  They are graph
observations: they do not include an ordered primitive trace witness.

This is a modeling boundary.  A concrete implementation may expose more
low-level detail, including timing or transport artifacts.  Such detail
belongs in a runtime machine and in a refinement or analysis assumption,
not in the language-independent event graph.  The event graph records
protocol events whose observations are part of the game description.

\subsection{Checkpoint batches}
\label{sec:mach:checkpoints}

Primitive execution can be batched.  A batch is ordered operational
evidence that a destination checkpoint is reachable from a source
checkpoint by available primitive events.  Strategic histories do not
store that ordered witness.  A checkpoint presentation instead specifies
which source/destination checkpoint pairs are allowed and proves that
each allowed checkpoint is realizable by some primitive batch.

The primitive downset checkpoint model is the canonical operational
presentation used for progress theorems.  The frontier strategic
presentation is more structured: it closes ready internal work, exposes
ready commit frontiers, applies selected player commits, and then closes
newly ready internal work before the next checkpoint observation.

A checkpoint should be read as a game-facing observation point, not as a
claim that the runtime literally pauses there.  It is the level at which
strategies are allowed to react.

\leancorr{The generic carrier is \TT{Vegas.Machine.Basic}; traces are in
\TT{Vegas.Machine.Trace}.  The primitive adapter is
\TT{Vegas.EventGraph.ToMachine}.  Batch and checkpoint definitions are
in \TT{Vegas.EventGraph.Batch} and
\TT{Vegas.Core.ToEventGraph.Checkpoint}.  The finite reachable-state
primitive machine is \TT{Vegas.Core.ToEventGraph.FiniteMachine}.}
